Une rame de métro part du repos et accélère uniformément à \qty{\v{a1}}{m/s^2} pendant \qty{\v{t1}}{s}. Elle roule ensuite à vitesse constante pendant \qty{\v{t2}}{s}, puis freine uniformément à \qty{\v{a3}}{m/s^2} jusqu'à l'arrêt à la station suivante.

\begin{enumerate}[label=\alph*)]
    \item Quelle est la vitesse maximale de la rame ?
    \item Quelle est la distance entre les deux stations ?
    \item Quelle est la vitesse moyenne de la rame sur tout le trajet ?
    \item Tracez l'allure du graphique de la vitesse en fonction du temps.
\end{enumerate}
